\section{Methods}
\label{sec:methods}

\subsection{Models}
The RQ1 continuation diagnostics evaluate seven pretrained, non-instruction-tuned autoregressive language models: Gemma 2 2B and 9B, Qwen3 4B and 8B, Llama 3 8B, Llama 3.2 3B, and OLMo 3 7B. The direct prompted judgments and RQ2--RQ4 use the same set except OLMo 3 7B. For representation analyses, we use the residual stream at the final prompt token. Computation uses bfloat16 model weights; extracted activations and fitted analysis objects are stored with fixed dataset and checkpoint hashes.

\subsection{Data}
The main causal training and development sets contain 2,000 rows (1,000 critical--control pairs) and 400 rows (200 pairs), respectively. The held-out causal evaluation set contains 300 rows (150 pairs). Within a pair, the critical item supports the target event as a causal consequence of the stated factor, whereas the control preserves the event and discourse structure while adding a non-explanatory factor. The revised training and development files are \texttt{rc\_modifier\_train\_set.txt} and \texttt{rc\_modifier\_dev\_set.txt}; these historical filenames contain the revised independent-clause stimuli used in this analysis. Held-out items are distinct from training and development items.

For the temporal ablation, we use 800 training, 200 development, and 200 test pairs. Positive and negative temporal items contrast an event-consistent temporal relation with a matched non-temporal control. The causal and temporal data are fit independently; cross-domain evaluations apply a frozen probe or direction from one domain to the other.

\subsection{RQ1: behavioral eliciture diagnostics}
RQ1 uses the 300 held-out causal rows (150 matched critical--control pairs), equally divided between relative-clause and modifier constructions. Scores use raw, untempered next-token log probabilities. For the ignorance diagnostic \citep{pan-kehler-2025-pragmatic}, we remove the context's terminal period and append the exact continuation ``\texttt{, and I don't know why.}'' The row score is the negative summed log probability of all continuation tokens,
\begin{equation}
 I_i=-\log P(\text{\emph{and I don't know why}}\mid x_i),
\end{equation}
so a positive critical-minus-control gap indicates that professing ignorance is less compatible with the inferred explanation in the critical item.

For the explicit-happenstance diagnostic, let $E_i$ be the event and $F_i$ the factor. Let $B_i$ denote the reconstruction prefix ending in $E_i$ followed by \emph{because}, and let $S_i$ denote the matched prefix ending in $E_i$ followed by \emph{and it so happened that}. We score the identical factor tokens after these prefixes and compute
\begin{equation}
 H_i=\overline{\log P}(F_i\mid B_i)-\overline{\log P}(F_i\mid S_i).
\end{equation}
The mean over factor tokens prevents target length from determining the score. A positive critical-minus-control gap means that causal framing provides a larger relative advantage for explanatory than irrelevant factors. RQ1 does not use the separate \emph{because}--versus--bare-\emph{and} diagnostic.

For each measure, we report critical and control means, the matched critical-minus-control gap, a 95\% interval from 10,000 pair-bootstrap samples, pairwise accuracy with ties worth one half, row-level ROC-AUC, and a one-sided paired sign-flip test. No threshold or prompt is selected on the test set. The explicit-happenstance contrast was prespecified as a secondary control in the original behavioral protocol, so its interval and unadjusted sign-flip test are interpreted separately from that protocol's primary FDR family.

\subsection{RQ1: prompted behavioral measurements}
For Yes/No judgments, we score pooled log-probability mass over capitalization and leading-space variants of \emph{Yes} and \emph{No} and normalize within the two answer classes,
\begin{equation}
 p_{\mathrm{yes}} = \frac{P(Y)}{P(Y)+P(N)}.
\end{equation}
We use both ``Answer Yes or No'' and ``Answer No or Yes.'' For Likert judgments, we score all integer candidates $k\in\{1,\ldots,7\}$ and use their posterior expectation,
\begin{equation}
 r = \sum_{k=1}^{7} k\frac{\exp q_k}{\sum_{j=1}^{7}\exp q_j},
 \qquad \tilde r=\frac{r-1}{6}.
\end{equation}
We use both ascending (1-to-7) and reversed (7-to-1) anchor descriptions. A score strictly above 0.5 is classified as positive; ties are classified as negative. We report each presentation order separately and an item-level average across the two orders. The external ceiling is one frozen 1--7 Gemini 3.1 Pro Preview rating per unchanged test item, normalized as $(r-1)/6$.

\subsection{RQ2: probes, CAA, and transfer}
For each candidate layer $\ell$, we fit a class-balanced logistic regression probe to distinguish critical from control items. The probe score is its raw decision logit. Regularization and the layer are selected on the corresponding development set, with the earliest layer breaking ties.

The CAA direction is the mean paired activation difference,
\begin{equation}
 v_\ell = \frac{1}{|\mathcal P|}\sum_{i\in\mathcal P}
 \left(h_{i,\ell}^{\mathrm{critical}}-h_{i,\ell}^{\mathrm{control}}\right).
\end{equation}
An item is scored by its projection onto the unit direction,
$s_{i,\ell}=h_{i,\ell}^{\top}v_\ell/\lVert v_\ell\rVert_2$.
CAA layers maximize the mean of development ROC-AUC and paired accuracy. Probe and CAA thresholds are chosen on the source-domain development split by balanced accuracy and then frozen. We report ROC-AUC, balanced accuracy, and paired accuracy on causal and temporal tests, including causal-to-temporal and temporal-to-causal transfer.

\subsection{RQ3: correlation with graded judgments}
We correlate the frozen causal probe logits and unit-direction CAA projections with (i) order-specific and order-averaged Yes/No probabilities, (ii) probabilistic Likert scores, and (iii) frozen Gemini ratings. Primary analyses use row-level Pearson and Spearman correlations. Secondary analyses first compute the critical-minus-control difference for each pair and correlate those differences. We keep representational scores raw: Pearson correlation is unchanged by positive affine rescaling, and Spearman correlation depends only on ranks. Ninety-five percent confidence intervals use 1,000 bootstrap resamples clustered by pair ID.

\subsection{RQ4: predicting failure modes}
Let $c\in\{0,1\}$ be the decision of the previously trained, frozen causal probe, using a raw-logit threshold of zero, and let $b\in\{0,1\}$ be the model's prompted behavioral decision. Each item--prompt observation receives one of three labels. \emph{Suppression} has $c=1,b=0$; \emph{post-hoc attribution} has $c=0,b=1$; and probe--output agreements ($b=c$) are \emph{correct}. The terms \emph{correct positive} and \emph{correct negative} therefore mean positive and negative agreement with the causal probe, not necessarily agreement with dataset ground truth. Appendix~\ref{app:failure-grid} visualizes the four combinations. Prompt-order variants of the same item may receive different failure labels.

At every layer, a class-balanced multinomial logistic regression receives the standardized final-token residual activation concatenated with a one-hot prompt-variant indicator. It is trained on all four formats in the semantically independent \texttt{rc\_modifier\_train\_set2.txt}. We select the layer on \texttt{rc\_modifier\_dev\_set2.txt} by macro-F1, then balanced accuracy, with the earliest layer as the final tie-breaker. The selected probe is frozen for the 1,200 held-out observations (300 items under four formats). We emphasize balanced accuracy and macro-F1 because the correct class is often much larger than either error class.

\subsection{RQ4: vector construction and intervention}
We construct separate endpoint-balanced vectors for Yes/No and Likert. For suppression at layer $\ell$,
\begin{equation}
\begin{aligned}
 v_{\mathrm{sup},\ell}={}&
 \mathbb E[h_\ell\mid\mathrm{suppression}]\\
 &-\mathbb E[h_\ell\mid\mathrm{correct\ positive}],
\end{aligned}
\end{equation}
For post-hoc attribution, the matched contrast reverses the gold label but retains the same failure-minus-correct orientation:
\begin{equation}
\begin{aligned}
 v_{\mathrm{post},\ell}={}&
 \mathbb E[h_\ell\mid\mathrm{post\mbox{-}hoc\ attribution}]\\
 &-\mathbb E[h_\ell\mid\mathrm{correct\ negative}].
\end{aligned}
\end{equation}
Contrasts are computed within each prompt order and then averaged, giving the two orders equal weight. Vectors are raw and are not length-normalized.

Vector layers maximize development AUC for separating the target failure from its matched-correct class. For $\alpha\in\{0,0.5,1,2,3\}$, we intervene at the selected layer with
\begin{equation}
 h'_\ell=h_\ell-\alpha v_{f,\ell}.
\end{equation}
Magnitude is selected on development data by balanced accuracy minus the rate at which initially correct examples become incorrect. At test time, the primary intervention is gated by the frozen failure probe; oracle gating is retained as a diagnostic. We report balanced accuracy, damage rate, net corrections, and the signed change in the Yes--No or Likert margin.
